\documentclass[11pt,a4paper]{article}
\usepackage[T1]{fontenc}
\usepackage[utf8]{inputenc}
\usepackage{amsmath,amssymb,graphicx,hyperref,longtable,booktabs,array}
\usepackage[margin=2.2cm]{geometry}
\providecommand{\tightlist}{\setlength{\itemsep}{0pt}\setlength{\parskip}{0pt}}
\setlength{\parindent}{0pt}\setlength{\parskip}{6pt}

\begin{document}
\hypertarget{piece-3-evaluation-harness-and-baselines-30-oct-2026}{%
\section{Piece 3: evaluation harness and baselines (30 Oct
2026)}\label{piece-3-evaluation-harness-and-baselines-30-oct-2026}}

Code: \texttt{code/harness.py}, \texttt{code/test\_harness.py},
\texttt{code/run\_baselines.py}, \texttt{code/plot\_baselines.py} (needs
\texttt{sensor.py} from Piece 2). Written by Claude (AI assistant); the
team must read and sign off before claiming it. Assumptions: headline
noise sd 8 mg/dL (2 and 5 as sensitivity), prototype meal sizes
unchanged, 5-min sampling, 5\% dropouts.

\hypertarget{protocol-implemented}{%
\subsection{Protocol implemented}\label{protocol-implemented}}

\begin{itemize}
\tightlist
\item
  27 usable virtual patients (same exclusions as the prototype).
  Subject-wise 5-fold CV (folds of 6, 6, 5, 5, 5; each has adolescents
  and adults). Manifest saved in \texttt{results/split\_manifest.json}.
\item
  Hyper-parameters are chosen by pooled RMSE on the four non-test folds
  only; test subjects are evaluated once. The baselines have no fitted
  parameters, so there is no separate train/validation split yet.
\item
  Estimators: B0 raw (last value carried over dropouts), B1 fixed time
  shift (non-causal), B2 causal exponential smoothing, B3 causal naive
  inversion with tau = 10, B3 oracle (told the true tau; a ceiling, not
  a fair competitor).
\item
  Metrics per subject: RMSE, MAD, bias, MARD (plasma \textgreater= 40),
  RMSE on rapid-change samples (\textbar true rate\textbar{}
  \textgreater= 1 mg/dL/min, at least 20 samples), rate-of-change RMSE
  (15-min window), residual lag (best shift), hypoglycaemia detection
  (threshold 70 mg/dL, events merged if gaps \textless{} 30 min,
  detection window 30 min before onset to end, false alarms per day).
  Summaries: median over subjects with cluster bootstrap 95\% CI (2000
  resamples); paired differences versus raw with Wilcoxon test and Holm
  correction; pooled-ratio bootstrap for hypoglycaemia delay and false
  alarms.
\end{itemize}

\hypertarget{tests-16-of-16-pass}{%
\subsection{Tests: 16 of 16 pass}\label{tests-16-of-16-pass}}

Hand-checked metrics (constant +10 error gives RMSE, MAD, bias, MARD all
10); perfect estimator; a 10-min lagged estimate gives residual lag 10
and hypo delay 10; false alarm counted; bootstrap and Holm helpers;
folds disjoint and complete; causality test (changing future samples
leaves past output of B0, B2, B3 unchanged, and B1 correctly fails it);
every subject evaluated once; noise-free sanity (oracle inverse RMSE
1.68 vs raw 7.29 mg/dL).

\hypertarget{results-at-noise-sd-8-mgdl-median-over-27-subjects-95-ci}{%
\subsection{Results at noise sd 8 mg/dL (median over 27 subjects, 95\%
CI)}\label{results-at-noise-sd-8-mgdl-median-over-27-subjects-95-ci}}

RMSE on rapid-change samples, mg/dL:

\begin{longtable}[]{@{}lllll@{}}
\toprule\noalign{}
Scenario & B0 raw & B1 shift (non-causal) & B3 inverse (causal, tau 10)
& B3 true tau (oracle) \\
\midrule\noalign{}
\endhead
\bottomrule\noalign{}
\endlastfoot
S0 matched & 21.5 {[}19.2, 24.7{]} & 11.3 {[}10.0, 12.6{]} & 17.9
{[}15.1, 20.2{]} & 17.7 {[}15.1, 20.0{]} \\
S1 sensor chain & 27.7 {[}23.8, 33.2{]} & 11.2 {[}9.9, 13.1{]} & 23.6
{[}19.9, 26.9{]} & 22.9 {[}18.7, 26.6{]} \\
S5 full & 31.3 {[}23.0, 32.9{]} & 12.4 {[}11.4, 14.0{]} & 26.3 {[}21.7,
27.3{]} & 25.5 {[}21.6, 28.5{]} \\
\end{longtable}

Paired change versus raw (rapid RMSE, Holm-adjusted p): S0 shift -10.3
{[}-12.5, -7.3{]} (p \textless{} 0.001), B3 -3.0 {[}-4.6, -1.3{]} (p =
0.001). S1: shift -16.5, B3 -4.6. S5: shift -18.4, B3 -4.2. All
scenarios and noise levels are in
\texttt{results/baseline\_results.json}.

\hypertarget{findings}{%
\subsection{Findings}\label{findings}}

\begin{enumerate}
\def\labelenumi{\arabic{enumi}.}
\tightlist
\item
  The non-causal shift (advance by 8 to 13 min, chosen in CV) cuts
  rapid-change error by about half in every scenario. It uses future
  samples, so it is a retrospective ceiling, not a real-time method.
  Causal inversion recovers only 3 to 5 mg/dL of the 10 to 18 mg/dL
  available at noise sd 8.
\item
  Overall RMSE (all samples): causal inversion does not beat raw at
  noise sd 8 (S0: 11.3 vs 11.1 mg/dL), consistent with the prototype.
  Its gain is confined to rapid-change periods and costs more false
  alarms (S0 1.30 vs 0.93 per day; S5 1.86 vs 1.19).
\item
  B2 causal smoothing picked alpha = 1 (no smoothing) in every fold, so
  it equals raw: smoothing adds lag and never lowers error versus truth.
  Kept as a sanity row.
\item
  B3 with the true tau is not better than B3 with tau = 10 at noise sd
  8: the local-slope estimate over a 25 to 30 min window leaves about 5
  min residual lag and noise limits the gain.
\item
  Rapid-change share varies from 0.0\% to 27.7\% of samples per patient
  (median 8.6\%); results are driven by the patients who move.
\item
  No problematic grid-edge picks: B2 hits alpha = 1, a hard physical
  limit; B1 and B3 optima are interior.
\item
  Hypoglycaemia: 57 to 58 events pooled across test subjects. Negative
  delays mean the estimate crossed 70 mg/dL before truth, mostly
  noise-driven; interpret delay together with false alarms per day.
\end{enumerate}
\end{document}
